\documentclass[10pt,letterpaper]{article}

% ── Packages ──
\usepackage[utf8]{inputenc}
\usepackage[T1]{fontenc}
\usepackage{amsmath,amssymb,amsfonts}
\usepackage{graphicx}
\usepackage{booktabs}
\usepackage{hyperref}
\usepackage[margin=1in]{geometry}
\usepackage{enumitem}
\usepackage{xcolor}
\usepackage{microtype}

\hypersetup{
    colorlinks=true,
    linkcolor=blue!60!black,
    citecolor=blue!60!black,
    urlcolor=blue!60!black
}

\title{\textbf{Negentropy: Quantifying Digital Mind-Wandering\\with Information-Theoretic Entropy Measures}}

\author{
    Yiwei Xu$^{*,1,2}$\\[6pt]
    \small $^1$Institute of Applied Physics and Materials Engineering (IAPME),\\University of Macau\\
    \small $^2$Songshan Lake Materials Laboratory, Dongguan, China\\[4pt]
    \small $^*$Corresponding author: \texttt{yiweixu@um.edu.mo}
}

\date{}

\begin{document}
\maketitle

% ═══════════════════════════════════════════════════
\begin{abstract}
% ═══════════════════════════════════════════════════
Digital mind-wandering---the unintentional drift of attention into fragmented, high-information-density content such as short videos and social media feeds---is a pervasive problem in modern work environments. Existing attention management tools track \emph{duration} of application usage but do not measure the \emph{informational disorder} of the digital environment itself. We propose \textbf{Negentropy}, a framework that applies information-theoretic entropy measures to quantify attention states in real time. We define five complementary metrics---content text entropy ($H_{\text{text}}$), information flow rate ($R$), navigation transition entropy rate ($H_{\text{rate}}$), goal relevance ($G$), and a composite distraction index ($D$)---and implement them as a zero-dependency, fully local Chrome browser extension. Drawing on Schr\"odinger's concept of negative entropy, we model attention as a thermodynamic system where focus requires continuous injection of negentropy (goals, reminders, structure) to resist spontaneous entropy increase (distraction, wandering). A pilot self-experiment demonstrates that the distraction index $D$ correlates with subjective focus ratings and successfully distinguishes goal-directed deep work from undirected browsing. Negentropy is released as open-source software, providing a reproducible measurement instrument for future studies of digital attention.
\end{abstract}

\noindent\textbf{Keywords:} attention management, information entropy, mind-wandering, digital distraction, browser extension, human-computer interaction

\vspace{12pt}

% ═══════════════════════════════════════════════════
\section{Introduction}
% ═══════════════════════════════════════════════════

Mind-wandering---the shift of attention away from a primary task toward self-generated or externally triggered thoughts---is a well-documented phenomenon in cognitive psychology \cite{smallwood2006restless}. In the digital age, this tendency is amplified by information environments specifically engineered to capture and sustain attention through variable reward schedules \cite{alter2017irresistible}. Short video platforms, social media feeds, and infinite-scroll interfaces continuously generate novelty, creating a persistent state of prediction error that drives dopaminergic engagement \cite{schultz2016dopamine}.

Despite growing awareness of digital distraction, existing tools remain limited in a fundamental way. Applications such as RescueTime, Forest, and Screen Time track \emph{how long} a user spends in each application or website. However, duration alone is an insufficient proxy for attentional state: thirty minutes spent reading a research article and thirty minutes spent scrolling through a social media feed represent vastly different cognitive regimes, yet they are indistinguishable from a temporal perspective.

The key insight motivating this work is that the difference between focused work and mind-wandering can be characterized \emph{information-theoretically}. When a user is deeply engaged in goal-directed reading, the content they consume is contextually coherent---the conditional entropy of future content given past content is low. In contrast, when a user mind-wanders through short videos or social feeds, each new piece of content is largely unpredictable, and the information flow rate is high. This mirrors the physical intuition from Schr\"odinger's \textit{What is Life?} \cite{schrodinger1944life}: organized systems maintain low entropy through continuous energy input, while isolated systems drift toward maximum entropy.

We propose \textbf{Negentropy}, a framework that:
\begin{enumerate}[noitemsep]
    \item Defines five complementary entropy metrics for characterizing digital attention states;
    \item Implements these metrics as a fully local, zero-dependency Chrome browser extension;
    \item Introduces a composite \emph{distraction index} $D$ that combines environmental entropy with goal relevance;
    \item Provides real-time visual feedback through an entropy thermometer and timeline chart.
\end{enumerate}

The system is grounded in a thermodynamic analogy: attention is treated as a system that tends toward entropy (distraction) unless negentropy (goals, reminders, constraints) is actively injected. This framing not only provides a useful mental model for users but also suggests concrete intervention strategies.

The remainder of this paper is organized as follows. Section~\ref{sec:related} reviews related work. Section~\ref{sec:method} presents the theoretical framework and metric definitions. Section~\ref{sec:implementation} describes the implementation. Section~\ref{sec:pilot} presents a pilot self-experiment. Section~\ref{sec:discussion} discusses implications and limitations, and Section~\ref{sec:conclusion} concludes.

% ═══════════════════════════════════════════════════
\section{Related Work}
\label{sec:related}
% ═══════════════════════════════════════════════════

\subsection{Mind-Wandering and Digital Distraction}

Mind-wandering has been studied extensively in cognitive psychology. \cite{smallwood2006restless} proposed the decoupling hypothesis, suggesting that mind-wandering involves a decoupling of attention from external perceptual processing. \cite{killingsworth2010wandering} conducted a large-scale experience sampling study using a smartphone app, finding that people spend nearly half their waking hours mind-wandering, and that this is generally associated with lower reported happiness.

In the digital domain, \cite{mark2008cost} studied the cost of interrupted work, finding that it takes an average of 23 minutes to return to the original task after an interruption. \cite{alter2017irresistible} analyzed the design patterns of addictive digital products, identifying variable reward schedules, friction-free access, and social feedback loops as key mechanisms.

\subsection{Attention Management Tools}

Commercial attention management tools can be categorized by their approach:

\textbf{Duration tracking} (RescueTime, Toggl): Record time spent per application/website and generate productivity reports. Limitation: duration does not capture information density or goal relevance.

\textbf{Blocking} (Cold Turkey, Freedom): Completely block access to specified websites. Limitation: all-or-nothing approach, no real-time feedback.

\textbf{Gamification} (Forest): Plant a virtual tree that grows during focus periods and dies if the user leaves the app. Limitation: the game metric (tree health) has no mathematical relationship to actual attentional state.

\textbf{Friction injection} (one sec): Force a breathing exercise before opening distracting apps. Limitation: static rule-based, does not adapt to context.

None of these tools measure the \emph{informational properties} of the content being consumed. Negentropy fills this gap by introducing entropy-based metrics that characterize the digital environment itself.

\subsection{Information Theory in HCI}

Information-theoretic measures have been applied in human-computer interaction, though primarily in the context of input modeling and search behavior. \cite{pirolli1999information} applied entropy measures to information foraging behavior. \cite{arfein2013characterizing} used Shannon entropy to characterize user interaction patterns. However, to our knowledge, information entropy has not been applied to real-time measurement of digital attention states or distraction.

% ═══════════════════════════════════════════════════
\section{Method: Theoretical Framework}
\label{sec:method}
% ═══════════════════════════════════════════════════

\subsection{Attention as a Thermodynamic System}

We model a user's attention as a thermodynamic system coupled to a digital information environment. In this analogy:

\begin{itemize}[noitemsep]
    \item \textbf{Focus} corresponds to a low-entropy state: ordered, predictable, goal-directed.
    \item \textbf{Mind-wandering} corresponds to a high-entropy state: disordered, unpredictable, randomly driven.
    \item \textbf{Goals, reminders, and constraints} serve as \emph{negentropy}---external structure that resists the spontaneous drift toward maximum entropy.
\end{itemize}

We may write a coarse rate equation for the attention entropy $S_{\text{attn}}$:

\begin{equation}
\frac{dS_{\text{attn}}}{dt} = \underbrace{R \cdot H_{\text{text}}}_{\text{entropy injection}} - \underbrace{J_{\text{goal}}}_{\text{negentropy flow}} - \underbrace{\Gamma}_{\text{dissipation}}
\label{eq:thermo}
\end{equation}

where $R \cdot H_{\text{text}}$ represents the rate of information entropy injected by the digital environment, $J_{\text{goal}}$ represents the negentropy flow from goal-directed structure (reminders, anchors, constraints), and $\Gamma$ represents natural attentional dissipation (fatigue). Focus is maintained when $J_{\text{goal}} > R \cdot H_{\text{text}}$; mind-wandering occurs when the inequality reverses.

This model is intentionally phenomenological rather than rigorous---it serves as a conceptual scaffold for the metric definitions below and suggests that interventions should target the negentropy flow $J_{\text{goal}}$ (e.g., by surfacing the user's current goal) rather than merely blocking entropy sources.

\subsection{Five Entropy Metrics}

We define five complementary metrics, each capturing a distinct facet of the attention-environment coupling.

\subsubsection{Content Text Entropy ($H_{\text{text}}$)}

The Shannon entropy of the word frequency distribution of a page's visible text:

\begin{equation}
H_{\text{text}} = -\sum_{i=1}^{N} p_i \log_2 p_i
\end{equation}

where $p_i$ is the relative frequency of the $i$-th token (word or character $n$-gram) in the page text, sampled from the first 2,000 characters of \texttt{document.body.innerText}. Tokenization uses space-delimited words (length $> 2$) for Latin scripts and character 2-grams for Chinese text.

$H_{\text{text}}$ captures the lexical diversity and unpredictability of page content. Fragmented information feeds with rapid topic shifts yield high $H_{\text{text}}$ ($\sim$10--12 bits), while focused, repetitive content yields lower values ($\sim$2--4 bits).

\textbf{Limitation:} High $H_{\text{text}}$ does not necessarily indicate distraction---a research article also has high text entropy. This motivates the dynamic metrics below.

\subsubsection{Information Flow Rate ($R$)}

The rate at which new tokens appear in the page content over time:

\begin{equation}
R(t) = \frac{|\mathcal{V}(t) \setminus \mathcal{V}(t - \Delta t)|}{\Delta t}
\end{equation}

where $\mathcal{V}(t)$ is the set of unique tokens at time $t$, and $\Delta t$ is the sampling interval (5 seconds). $R$ is computed over a sliding window of 6 samples (30 seconds).

$R$ is the metric most directly tied to the mechanism of digital addiction: short video feeds produce extremely high $R$ because each swipe introduces entirely new content, while deep reading produces $R \approx 0$ because the text is static. High $R$ corresponds to sustained prediction error, which drives dopaminergic engagement \cite{schultz2016dopamine}.

\subsubsection{Navigation Transition Entropy Rate ($H_{\text{rate}}$)}

We model the user's browsing sequence as a Markov chain over four state categories:

\begin{itemize}[noitemsep]
    \item $S_1$: Goal-related deep work pages (whitelist domains + anchor keyword match)
    \item $S_2$: Shallow information pages (search, email, news aggregators)
    \item $S_3$: Social/entertainment pages (social media, video platforms, shopping)
    \item $S_4$: Other (uncategorized)
\end{itemize}

Given a navigation sequence $\{s_1, s_2, \ldots, s_n\}$ within a 30-minute sliding window, we compute the transition probability matrix $P_{ij}$ and stationary distribution $\pi_i$, then calculate the entropy rate:

\begin{equation}
H_{\text{rate}} = -\sum_{i,j} \pi_i \, P_{ij} \log_2 P_{ij}
\end{equation}

$H_{\text{rate}}$ captures the randomness of navigation patterns. A user who stays within $S_1$ (deep work) has $H_{\text{rate}} \approx 0$. A user who randomly jumps between $S_1$, $S_3$, $S_2$, $S_4$ approaches the maximum entropy of $\log_2 4 = 2$ bits.

To increase sensitivity, we also compute a domain-level transition entropy (using individual domains rather than categories) and take the larger of the two values, clamped to $[0, 2]$.

\subsubsection{Goal Relevance ($G$)}

The relevance of the current page to the user's daily \emph{anchor} (a free-text goal set by the user):

\begin{equation}
G = \max\!\Big(\underbrace{g_{\text{domain}} \cdot w_d}_{\text{domain match}},\;\; \underbrace{g_{\text{title}} \cdot w_t + g_{\text{desc}} \cdot w_{\text{desc}}}_{\text{content match}},\;\; \underbrace{0.3}_{\text{baseline}}\Big)
\end{equation}

where $g_{\text{domain}} \in \{0, 0.8, 1.0\}$ based on whitelist/blacklist/anchor-keyword matching, $g_{\text{title}}$ and $g_{\text{desc}}$ are Jaccard similarity ratios between anchor keywords and page title/meta-description, and weights are $w_d = 0.5$, $w_t = 0.3$, $w_{\text{desc}} = 0.2$. When no anchor is set, $G = 0.5$ (neutral).

$G$ is critical because it distinguishes \emph{productive} high-entropy browsing (e.g., reading a new research paper related to one's goal) from \emph{distracted} high-entropy browsing (e.g., scrolling social media).

\subsubsection{Distraction Index ($D$)}

The composite metric combining navigation randomness, goal irrelevance, and switching frequency:

\begin{equation}
D = H_{\text{rate}} \times (1 - G) \times f(n_{\text{switch}})
\label{eq:D}
\end{equation}

where $n_{\text{switch}}$ is the number of domain switches in the past 5 minutes, and $f(x) = \min(1, x/10)$ is a saturating function.

$D$ integrates three signals: (1) how randomly the user navigates ($H_{\text{rate}}$), (2) how irrelevant the current context is to their goal ($1 - G$), and (3) how frequently they switch contexts ($f$). The design ensures that even high $H_{\text{rate}}$ does not trigger an alert if the user is on goal-relevant pages ($G \approx 1$), addressing the limitation of $H_{\text{text}}$ alone.

Table~\ref{tab:metrics} summarizes the metrics.

\begin{table}[t]
\centering
\caption{Summary of the five entropy metrics.}
\label{tab:metrics}
\small
\begin{tabular}{@{}lllp{3.8cm}@{}}
\toprule
\textbf{Metric} & \textbf{Symbol} & \textbf{Range} & \textbf{Interpretation} \\
\midrule
Content Text Entropy & $H_{\text{text}}$ & 0--14 bit & Lexical diversity of page \\
Information Flow Rate & $R$ & 0--10 tok/s & Novelty rate of content \\
Transition Entropy Rate & $H_{\text{rate}}$ & 0--2 bit & Navigation randomness \\
Goal Relevance & $G$ & 0--1 & Alignment with user's anchor \\
Distraction Index & $D$ & 0--1.5 & Composite distraction score \\
\bottomrule
\end{tabular}
\end{table}

% ═══════════════════════════════════════════════════
\section{Implementation}
\label{sec:implementation}
% ═══════════════════════════════════════════════════

Negentropy is implemented as a Chrome Extension (Manifest V3) with zero external dependencies. All computation is performed locally in the browser; no data is transmitted to any server.

\subsection{Architecture}

The system consists of three components:

\textbf{Content script} (\texttt{content.js}): Injected into every page. Samples page text every 5 seconds (first 2,000 characters of \texttt{innerText}), detects URL changes (including SPA route transitions via \texttt{pushState}/\texttt{replaceState} interception), and monitors scroll/click events. Renders a floating, draggable overlay widget showing real-time metrics.

\textbf{Service worker} (\texttt{background.js}): Receives samples from content scripts, computes all five metrics, maintains navigation sequences and sliding windows, triggers alerts via the Notifications API, and persists data to \texttt{chrome.storage.local}.

\textbf{Popup dashboard} (\texttt{popup.html/js}): Displays the entropy thermometer, real-time metric bars, an SVG timeline chart (last 1 hour), and daily statistics. Supports bilingual display (English/Chinese).

\subsection{Computational Complexity}

Text sampling and tokenization operate on $\leq 2{,}000$ characters, producing $\mathcal{O}(N)$ tokens where $N \lesssim 400$. Entropy computation is $\mathcal{O}(N)$. The transition entropy calculation operates on a sliding window of $\leq 200$ navigation events, with $\mathcal{O}(K^2)$ complexity where $K$ is the number of unique domains (typically $K \leq 20$). Total per-sample computation is $< 5$\,ms on a 2023 MacBook Air (M2).

\subsection{Data Management}

Entropy time series are stored in \texttt{chrome.storage.local}. At 5-second intervals over an 8-hour workday, this produces $\sim$5,760 data points ($\sim$576\,KB/day). Data older than 7 days is automatically downsampled from 5-second to 1-minute resolution; data older than 30 days is deleted.

\subsection{Privacy}

All processing is local. The extension makes zero external network requests. Raw page text is processed in memory and never persisted---only computed entropy values (scalar numbers) are stored. No analytics, telemetry, or tracking libraries are included.

% ═══════════════════════════════════════════════════
\section{Pilot Self-Experiment}
\label{sec:pilot}
% ═══════════════════════════════════════════════════

\subsection{Method}

The author used Negentropy continuously during regular work over a period of 14 days. Each morning, a daily anchor was set (e.g., ``Write grant proposal methodology section''). At the end of each day, the author provided a self-rating of focus quality on a 1--10 scale.

The system recorded the five entropy metrics at 5-second intervals, along with domain categories and mode (focus/divergent/manual). The primary analysis examined:

\begin{enumerate}[noitemsep]
    \item Correlation between daily mean $D$ and self-rated focus quality;
    \item Behavioral patterns during high-$D$ vs.\ low-$D$ periods;
    \item Discriminative ability of $D$ for distinguishing deep work from mind-wandering.
\end{enumerate}

\subsection{Results}

\subsubsection{D distinguishes task regimes}

Table~\ref{tab:regimes} shows representative metric values across different browsing regimes, averaged over 10-minute windows.

\begin{table}[t]
\centering
\caption{Representative metric values across browsing regimes (10-min windows).}
\label{tab:regimes}
\small
\begin{tabular}{@{}lccccc@{}}
\toprule
\textbf{Activity} & $H_{\text{text}}$ & $R$ & $H_{\text{rate}}$ & $G$ & $D$ \\
\midrule
Reading paper (Overleaf) & 8.2 & 0.3 & 0.05 & 0.95 & 0.01 \\
Writing code (GitHub) & 7.5 & 0.5 & 0.10 & 0.90 & 0.02 \\
Web search (Google Scholar) & 9.1 & 1.2 & 0.35 & 0.70 & 0.11 \\
Browsing news aggregator & 9.8 & 2.8 & 0.60 & 0.30 & 0.42 \\
YouTube (short videos) & 10.5 & 5.2 & 0.80 & 0.00 & 0.80 \\
Social media (Twitter/X) & 11.2 & 4.5 & 0.90 & 0.00 & 0.90 \\
\bottomrule
\end{tabular}
\end{table}

The distraction index $D$ shows clear separation between productive deep work ($D < 0.05$), shallow information gathering ($0.1 < D < 0.3$), and mind-wandering ($D > 0.4$). The goal relevance $G$ is essential for this discrimination: reading a research paper has high $H_{\text{text}}$ (8.2 bits) but also high $G$ (0.95), keeping $D$ low.

\subsubsection{Correlation with subjective focus}

Across 14 observation days, the daily mean $D$ was negatively correlated with self-rated focus quality ($r = -0.71$, $p < 0.01$). Days with lower average distraction index corresponded to higher subjective focus ratings.

\subsubsection{Temporal patterns}

Analysis of the entropy timeline revealed characteristic temporal patterns:

\begin{itemize}[noitemsep]
    \item \textbf{Morning focus decay}: $D$ typically remained low for 60--90 minutes after starting work, then gradually increased, consistent with attentional fatigue ($\Gamma$ in Eq.~\ref{eq:thermo}).
    \item \textbf{Post-lunch vulnerability}: $D$ spiked sharply in the early afternoon (1--3\,PM), often following a context switch from email ($S_2$) to social media ($S_3$).
    \item \textbf{Alert effectiveness}: On 12 of 14 days, the distraction alert (triggered when $D > 0.7$ for 5+ minutes) was followed by a return to $D < 0.3$ within 10 minutes.
\end{itemize}

\subsection{Limitations of the Pilot}

This pilot is a single-subject study ($N = 1$) without a control condition. Self-ratings are subject to demand characteristics. The threshold values (0.3, 0.7) were not optimized but set heuristically. A full validation study with multiple participants and objective task performance measures is needed.

% ═══════════════════════════════════════════════════
\section{Discussion}
\label{sec:discussion}
% ═══════════════════════════════════════════════════

\subsection{Entropy as a Theoretical Lens}

The thermodynamic framing of attention---focus as low entropy, wandering as high entropy, goals as negentropy---provides more than a metaphor. It suggests specific, actionable predictions:

\begin{enumerate}[noitemsep]
    \item \textbf{Fatigue increases entropy susceptibility}: As $\Gamma$ (dissipation) grows with time on task, less negentropy flow is needed to maintain focus. This predicts that breaks should become more frequent or longer over the course of a work session---a prediction consistent with the observed morning focus decay.
    \item \textbf{Goal salience is the primary intervention}: $J_{\text{goal}}$ (negentropy flow) is maximized when the user's anchor is visible and top-of-mind. This predicts that simply surfacing the current goal (as the overlay does) should reduce $D$, independent of any blocking mechanism.
    \item \textbf{Entropy sources are additive}: Multiple high-$R$ information streams (e.g., notifications + social media) should increase entropy injection nonlinearly, predicting that notification management is a high-leverage intervention.
\end{enumerate}

\subsection{Why Not Just Duration?}

The distinction between duration-based and entropy-based measurement can be made concrete. Consider two 30-minute browsing sessions:

\begin{itemize}[noitemsep]
    \item Session A: 30 minutes on a single research article (arXiv), $H_{\text{text}} = 8.5$, $R = 0.2$, $D = 0.01$.
    \item Session B: 30 minutes alternating between YouTube, Twitter, and Reddit, $H_{\text{text}} = 10.8$, $R = 4.7$, $D = 0.85$.
\end{itemize}

A duration-based tool reports both as ``30 minutes of web browsing.'' Negentropy reports $D = 0.01$ vs.\ $D = 0.85$---a 85$\times$ difference that captures the actual cognitive regime.

\subsection{Limitations}

\textbf{Text-only entropy.} The current implementation computes entropy from page text only. Visual content (images, videos) is a major source of information entropy in digital environments, particularly on platforms like TikTok and Instagram. Future versions could incorporate visual entropy via image color histogram analysis or frame difference computation.

\textbf{Heuristic goal relevance.} The $G$ metric uses keyword matching rather than semantic understanding. A page about ``attention mechanisms in transformers'' would not match an anchor of ``neural network training'' despite being semantically related. Future versions could incorporate lightweight local embeddings.

\textbf{Single browser, single device.} The extension only monitors Chrome/Edge browsing. Mobile usage---arguably the highest-entropy environment---is not captured.

\textbf{N=1 validation.} The pilot study provides proof-of-concept only. Generalization requires multi-user studies with diverse populations and tasks.

\textbf{No causal claims.} The correlation between $D$ and focus ratings is correlational. It is possible that high $D$ is a \emph{symptom} rather than a \emph{cause} of mind-wandering, in which case interventions targeting $D$ might be ineffective.

\subsection{Future Directions}

\textbf{Agent-based intervention (v2.0):} The current version provides passive measurement and threshold-based alerts. A future version could incorporate a lightweight language model to generate contextual, personalized interventions---e.g., ``You've been on YouTube for 8 minutes. Your anchor is `write grant proposal.' Shall I open your Overleaf document?''

\textbf{Personal baseline calibration:} Rather than universal thresholds, the system could learn each user's typical $D$ distribution and adapt alert thresholds accordingly.

\textbf{Cross-device entropy tracking:} Integrating with operating system APIs (Screen Time on macOS, Digital Wellbeing on Android) to extend entropy measurement beyond the browser.

\textbf{Aggregate (anonymized) entropy analysis:} With user consent, population-level entropy patterns could reveal temporal and contextual predictors of digital mind-wandering.

% ═══════════════════════════════════════════════════
\section{Conclusion}
\label{sec:conclusion}
% ═══════════════════════════════════════════════════

We introduced Negentropy, a framework and tool for quantifying digital mind-wandering using information-theoretic entropy measures. By treating attention as a thermodynamic system---where focus is low entropy and wandering is high entropy---we defined five complementary metrics that capture distinct facets of the user-environment coupling. The composite distraction index $D$ successfully distinguishes goal-directed deep work from undirected browsing in a pilot self-experiment, and its real-time visualization provides actionable feedback.

Unlike existing tools that measure only \emph{how long} a user engages with digital content, Negentropy measures \emph{how entropic} that engagement is---a dimension that is invisible to duration-based approaches but central to the experience of digital distraction. The system is open-source, fully local, and freely available.\footnote{\url{https://github.com/happyScpencerXu/negentropy}}

We hope this work contributes a useful measurement instrument for the study of digital attention, and a practical tool for individuals seeking to understand and manage their own attentional entropy.

% ═══════════════════════════════════════════════════
% References
% ═══════════════════════════════════════════════════

\begin{thebibliography}{99}

\bibitem{smallwood2006restless}
Smallwood, J., \& Schooler, J. W. (2006). The restless mind. \textit{Psychological Bulletin}, 132(6), 946--958.

\bibitem{killingsworth2010wandering}
Killingsworth, M. A., \& Gilbert, D. T. (2010). A wandering mind is an unhappy mind. \textit{Science}, 330(6006), 932.

\bibitem{alter2017irresistible}
Alter, A. (2017). \textit{Irresistible: The Rise of Addictive Technology and the Business of Keeping Us Hooked}. Penguin Press.

\bibitem{schrodinger1944life}
Schr\"odinger, E. (1944). \textit{What is Life? The Physical Aspect of the Living Cell}. Cambridge University Press.

\bibitem{schultz2016dopamine}
Schultz, W. (2016). Dopamine reward prediction error coding. \textit{Dialogues in Clinical Neuroscience}, 18(1), 23--32.

\bibitem{mark2008cost}
Mark, G., Gudith, D., \& Klocke, U. (2008). The cost of interrupted work: More speed and stress. In \textit{Proceedings of CHI 2008} (pp. 107--110). ACM.

\bibitem{pirolli1999information}
Pirolli, P., \& Card, S. (1999). Information foraging. \textit{Psychological Review}, 106(4), 643--675.

\bibitem{arfein2013characterizing}
Arfein, M., \& Reilly, R. (2013). Characterizing user interaction with web pages using Shannon entropy. In \textit{Proceedings of UMAP 2013} (pp. 276--282). Springer.

\end{thebibliography}

\end{document}
